% Ad Atticum 16.1 --- headnote
% ad-atticum-16-01, 8 July 44 BC, in Puteolano

\textit{Cicero to Atticus, written at the Puteolan villa on
8 July 44 BC --- Perseus dateline \textit{Scr. in Puteolano
viii Id. Quint. a. 710 (44)}. This letter opens Book 16 of the
correspondence with Atticus, the last surviving book of the
collection; what follows runs from July through early November
44 BC, the final months in which Cicero's letters to Atticus
are preserved. Cicero has reached Puteoli on the Nones of July
on his way south toward embarkation for Greece. The first
indignation of the letter is the calendar: the month formerly
called \textit{Quinctilis} has just been renamed \textit{Iulius}
in Caesar's honour --- so Atticus's letter is dated ``the Nones
of July,'' and Cicero, writing to a Brutus who has lately killed
the man so commemorated, finds the new month-name unbearable
(``By Hercules, what a crew they are!''; ``Could anything be
more disgraceful than `July' for Brutus?''). His response is
the Greek tag \textit{et' e\=omen}, ``let it still be,'' an old
Homeric refrain --- leave it, drop it, let it pass.}

\textit{The body of the letter is a rapid sweep of the news.
The Buthrotian land-grant trouble (Caesar's veterans on
Atticus's Epirote estates) has flared into violence; Plancus
is reportedly moving fast; Pompey the younger may or may not be
in arms in Spain (Cicero takes the rumour of Sextus as serious
enough that, if true, ``we must serve as slaves, but without a
civil war''). Ventidius's threat he writes off as a panic. The
financial paragraph is unusually granular: 210,000 sesterces
have been arranged; young Marcus at Athens is to have his
allowance adjusted to 80,000 sesterces a year from the Kalends
of April, instead of being kept on the meager driblets Xeno has
been doling out. The closing paragraph is the most personal:
young Quintus, Cicero's nephew --- previously a great fraud ---
has now promised he will be ``a Cato,'' and asks his uncle to
stand surety for the reformation. Cicero will not vouch for
him; he sends Atticus the letter the boy has dictated and asks
him to judge for himself, but quietly warns that he himself is
unmoved. ``Heaven grant he does what he promises: the joy
would be common to us both. But I --- I will say no more.''}
